\def\ChapterCode{URO}
% \clearpage % TEMP temp clearpage
\chapter
[\CO{[URO]} Urologische Erkrankungen]
{Urologische Erkrankungen}
\ChapterIntro
{%
Erkrankungen des harnableitenden Systems.
}
{%
\begin{description}
  \item[Maintainer:] Sebastian Gabriel
  \item[Autoren:] Diverse
  \item[Reviewer:] Standard-Reviewprozess
  \item[Version:] {\VarDocVersionTypeName} ({\VarDocVersionTypeDescription})
  \item[SHA1:] {\DocGitCommit}
\end{description}

{\TxTeilkapitelAASS}
}


\clearpage
\section{Exkurs: \HeadingKeyword{Harnkatheter}}

\Z{\HeadingSub{\Lang{Sny.} Transurethraler Blasenkatheter}}

\TextSlide{\TxBeschreibung}
{Ein \I{Harnkatheter} ist ein harnableitendes System, bei dem ein
Schlauch (Katheter) über die Harnröhre in die Harnblase
eingeführt wird um Harn abzulassen. Mittels eines
aufblasbaren Ballons kann der Katheter in der Blase fixiert
werden, man spricht dann vom \I{Dauerkatheter}. Am
Harnkatheter kann ein \II{Harnbeutel} angeschlossen werden,
welcher den Urin sammelt und sich entleeren lässt. Der
Harnbeutel muss \EE{unter dem Patienten angebracht} sein, da
andernfalls der Urin zurückrinnen würde. Beim Umlagern
ist daher der zum Patienten führende \EE{Schlauch
abzuklemmen}. Der Katheter und das daran hängende System ist
eine mögliche \EE{Eintrittspforte für Keime}, daher ist auf einen
sauberen, kontaminationsfreien Umgang zu achten.
Wenn der Harnbeutel ausgelassen werden muss, so ist die
abgelassene Menge zu notieren, da bei manchen Patienten der
Flüssigkeitshaushalt bilanziert wird.

Ein Harnkatheter ist eine
medizinisch-pflegerische Routinemaßnahme und entsprechend
oft bei Patienten anzutreffen.}
{
\begin{itemize}
  \item Katheter leitet aus Harnblase Urin ab
  \item Harnbeutel -- \E{unterhalb} d. Patienten
  \item Schlauch abklemmen beim Umlagern
  \item Wenn ablassen, Menge notieren
  \item Infektionsgefahr
\end{itemize}
}
{}
{}
{}



\begin{PictureSeriesWide}{}{Bilderserie: Harnkatheter}
\PictureSeriesRow{3}
{./images/demaw/Skriptumaegf-img52-thirdwidth}
{Ein Harnkatheter \SourcePicture{DEMAW/Blacky}}
{./images/demaw/Skriptumaegf-img53-thirdwidth}
{Ein Harnbeutel \SourcePicture{DEMAW/Blacky}}
{./images/demaw/Skriptumaegf-img59-thirdwidth}
{Einführen eines Harnkatheters
unter sterilen Bedingungen \SourcePicture{DEMAW/Blacky}}
{empty}{}
\end{PictureSeriesWide}


\section{\HeadingKeyword{Nierenkolik}}
\label{topic:nierenkolik}%
\index{Kolik!Nieren-}%
\index{Nierenkolik}%


\TextSlide{{\TxBeschreibung} und Ursachen}{%
\DefinitionInPara{Eine \T{Nierenkolik} ist
gekennzeichnet durch kolikartige Schmerzen aufgrund einer  
Blockade des Harnleiters und dem daraus folgenden Harnstau.}
Vgl. Kolik: \FullPrettyRef{topic:kolik}.
Sie ist eine äußerst schmerzhafte
akute Erkrankung und entsteht meist dadurch, dass
Nierensteine den Harnleiter blockieren bzw. verstopfen.
Dadurch wird der Harn gestaut und der Harnleiter gedehnt
dadruch kommt es zu einem Dehnungsschmerz.
}{
\begin{itemize}
  \item Blockade des Harnleiters $\rightarrow$ Harnstau
  \item Dehnung des Harnleiters $\rightarrow$ Kolikartiger Schmerz
  \item Nierenstein(e)
\end{itemize}
}{}{}{}

\TextSlide{\TxSymptome}{%
Die Folge sind massive,
kolikartige Schmerzen, in der Flanke bzw. im Unterbauch der
jeweiligen Seite, welcher auch in den Hoden bzw. in die
Schamlippen ausstrahlen kann. Im Gegensatz zu anderen
abdominellen Erkrankungen ist der Patient extrem unruhig und
agitiert (\Speech{\enquote{Der Stein wandert, der Schmerz
wandert, der
Patient wandert}}).
Oft läuft er gestikulierend und wehklagend durch die Gegend.
Mitunter ist eine entsprechende Vorerkrankung
\E{(Nierensteine)} bereits bekannt und kann erfragt werden.
Unter Umständen kann Blut im Harn sein.

\bigskip


}{
\begin{itemize}
    \item Rücken-/Flankenschmerzen
    \item Evtl. Unterbauchschmerzen bei  Stein-/Sandabgang
    \item \Speech{\enquote{Der Stein wandert, der Schmerz wandert, der Patient
    wandert}}
    \item Blut im Harn (Hämaturie)
\end{itemize}
}{}{}{}

\SlideCompensationNoParskip{1}

% M %%%%%%%%%%%%%%%%%%%%%%%%%%%%%%%%%%%%%%%%%%%%%%%%%%% M %
% Nierenkolik
\ProcedurePrintDefault{MN23XX0C}


\section{\HeadingKeyword{Harnwegsinfekt} und Nierenbeckenentzüdung }
\label{topic:harnwegsinfekt}%
\label{topic:nierenbeckenentzuendung}%


\TextSlide{\TxBeschreibung}
{%
\DefinitionInPara{Als \T{Harnwegsinfekt} wird allgemein eine
Infektion der Harnwege bezeichnet.} Besonders gefährdet sind
Frauen, infolge der kürzeren Harnröhre und der
normalen bakteriellen Besiedelung der Scheide. Die
häufigsten Erreger sind, aufgrund der räumlichen
Nähe, klassische Darmkeime \Z{(E.
coli)}.
Wird ein Harnwegsinfekt nicht behandelt, kann die
Infektion aufsteigen und eine Nierenbeckenentzüdung
verursachen. \DefinitionInPara{Eine
\T{Nierenbeckenentzüdung} ist eine Entzündung des
Nierenbeckens. \Z{\Lang{Syn.} \I{Pyelonephritis}.}} 
Sie ist oft bakteriell bedingt 
und umfasst die oberen Harnwege, das Nierenbeckenkelchsystem
und u.\,U. auch das Nierengewebe und ist eine der häufigsten
Erkrankungen der Niere \cite{Pschy259}. 
Im schlimmsten Fall kann sich
die Infektion weiter ausbreiten und eine Sepsis
verursachen (\FullPrettyRef{topic:sepsis}).

Harnwegsinfektionen werden häufig in
Gesundheitseinrichtungen übertragen
(\enquote*{Spitalsinfektion}).
Problematisch ist hierbei, dass bei derartigen Infektionen
häufig multiresistente Keime (\enquote*{Spitalskeime},
gegen Antibiotika unempfindliche Bakterien,
\FullPrettyRef{topic:multiresistente-keime}) der Auslöser
sind, welche deutlich schwieriger zu behandeln sind. Ein
besonderes Risiko stellen hierbei harnableitende Systeme
(Harnkatheter, -beutel) da.
}
{%
\begin{itemize}
  \item Infektion der Harnwege
  \item Oft: Darmkeime
\end{itemize}
}
{}
{}
{}



\TextSlide{\TxSymptome}
{Der Patient klagt über Schmerzen beim Harnlassen, sowie
häufigen, aber nicht ergiebigen Harndrang. Es können
Schmerzen im Nierenlager oder der Blasengegend bestehen.
Allgemeinsymptome wie Fieber, Müdigkeit und
\enquote*{Krankheitsgefühl} sind möglich. 

Bei der akuten
Nierenbeckenentzündung leidet der Patient unter Fieber,
Flankenschmerzen, häufigen, aber nicht ergiebigen Harndrang
sowie Schmerzen beim urinieren. Ein chronischer Verlauf ist 
oft symptomarm.
}
{
\begin{itemize}
  \item Harnwegsinfekt:
  \begin{itemize}
    \item Schmerzen beim urinieren, evtl. im Nierenlager,
    Blasengegend
    \item Häufiger, unergiebiger Harndrang
    \item Fieber
  \end{itemize}
  \item Nierenbeckenentzüdung:
  \begin{itemize}
    \item Evtl. Flankenschmerz
  \end{itemize}
\end{itemize}
}
{}
{}
{}


\section{Akutes \HeadingKeyword{Harnverhalten}}
\label{topic:harnverhalten}%

\TextSlide{{\TxBeschreibung} und Ursachen}
{%
\DefinitionInPara{Ein \T{Akutes
Harnverhalten} ist Abflussbehinderung des Harns aus der
Harnblase, bei der verschiedene Ursachen möglich sind.}
Es ist besonders bei Männern im
\enquote{besten} oder fortgeschrittenem Alter häufig
anzutreffen. Hauptursache  dafür ist eine altersbedingte
Vergrößerung der
Prostata, die die Harnröhre abschnürt. Davon abgesehen
können auch alte Narben, Medikamente oder eine Störung der
Schließmuskel der Harnblase\footnote{zum Beispiel im Rahmen
eines Querschnitt-Syndroms nach Wirbelsäulenverletzungen} zu
Harnverhalten führen.

Je nach Füllungszustand der Harnblase kann das
Harnverhalten sehr schmerzhaft sein.
Die Maßnahmen beschränken sich auf einen zügigen Transport
an eine Abteilung für Urologie. Eine Notarzt-Nachforderung
ist kaum sinnvoll, da ein \Abk{NEF} oder \Abk{NAW} selten geeignetes
Material zur Harnableitung mitführt. 
}
{
\begin{itemize}
  \item Vergrößerung der \EE{Prostata}
  \item Lähmung (Querschnittlähmung)
  \item Schließmuskel-Krampf
  \item Verlegung der Harnröhre
  \item Medikamente
\end{itemize}
}{}{}{}

% M %%%%%%%%%%%%%%%%%%%%%%%%%%%%%%%%%%%%%%%%%%%%%%%%%%% M %
% Akutes Harnverhalten
\ProcedurePrintDefault{MR33XX1C}



\section{\HeadingKeyword{Niereninsuffizienz} und Nierenversagen}
\label{topic:niereninsuffizienz}
\label{topic:cni-therapie-dialyse}
\label{topic:dialysepatienten-gefahren}

\AuthNotes{Abkz. CNI}

Meistens kommt es im Rahmen anderer Krankheiten (Diabetes
mellitus, Hypertonie, erhöhte Blutfette \ldots) zu einer
massiven Schädigung der kleinen Nierengefäße und in weiterer
Folge zum Nierenversagen. Da die zuvor erwähnten Krankheiten
sehr häufig sind, ist in weiterer Folge die
Niereninsuffizienz ebenfalls häufig anzutreffen.

Als Therapie steht der Medizin anfänglich eine medikamentöse
Behandlung, später jedoch nur die Nierentransplantation oder
die \E{Blutwäsche} (\T{Dialyse}\index{Dialyse}) zur
Verfügung. Bei der klassischen Dialyse wird der Patient jeden
zweiten oder dritten Tag für einige Stunden an eine Maschine
angeschlossen, welche sein Blut reinigen kann. Die
\enquote*{Dialysefahrten}, also Transport zu einem oder von
einem Dialysezentrum, sind eine der Hauptaufgaben des
Krankentransportdienstes. Dabei ergeben sich folgende
Besonderheiten: \index{Dialyse-Patient!Gefahren}

\begin{itemize}
  \item Dialyse-Patienten haben häufig einen
  \EE{Dialyse-Shunt} (\FullPrettyRef{topic:shunt}).
  \item \EE{Der Ausfall der normalen Nierenfunktion hat Auswirkungen}
  auf den Wasserhaushalt, den Elektrolythaushalt und den
  Säure-Basen-Haushalt. Diese Funktionen der Niere sind zum
  Überleben elementar (Vitalfunktionen zweiter Ordnung).
  Dialysepatienten müssen daher sehr darauf achten nicht zu
  viel zu trinken. Weiters können \EE{Elektrolytstörungen}
  sogar \Ca{lebensbedrohliche Notfälle} nach sich ziehen --
  bis hin zum Herzstillstand!
  \item Die \EE{Dialyse} stellt einen schwerwiegenden
  Eingriff in den Wasser- und Elektrolythaushalt dar, dementsprechend
  kann es auch nach einer Dialyse zu den oben genannten
  \EE{Elektrolytstörungen} kommen!
\end{itemize}

\begin{Danger}
  \item Dialysepatientienten sind immer sorgsam zu
  überwachen. Lebensbedrohliche Notfälle sind aufgrund einer
  Elektrolytentgleisung, besonders unmittelbar vor und nach
  einer Dialyse, möglich!
\end{Danger}






% \subsection{Hodentorsion}
% 
% Verdrehung eine Hodens im Hodensack  ={\textgreater}
% 
% Abschnürung von Samenstrang und  Blutgefäßen
% Oft nach Eintritt in Pubertät, ev. {\textgreater} 20 J.
% 
% \paragraph{Symptome:}
% 
% {\mdseries
% akuter heftiger Schmerz}
% 
% {\mdseries
% Schwellung}
% 
% \paragraph{Therapie:}
% 
% hochlagern
% 
% umgehende urologisch-chirurgische  Intervention (ad URO!!!)




% \clearpage

% \section*{\protect\dsliterary{} Weiterführende Literatur}
% \begin{WidePar}\tiny
% \begin{multicols}{2}
% \Z{%
% \begin{labeling}{mmm}
%   \item[\citep{Renz-Polster2006}] \emph{\bibentry{Renz-Polster2006}} 
%   \item[\citep{SiewertChirurgie8}] \emph{\bibentry{SiewertChirurgie8}}
%   \item[\citep{ScholzNotfallmedizin2}] \emph{\bibentry{ScholzNotfallmedizin2}} 
%   \item[\citep{IntroClinEmergMed4}] \emph{\bibentry{IntroClinEmergMed4}} 
%   \item[\citep{Longmore2007}] \emph{\bibentry{Longmore2007}} 
%   \item[\citep{Fauci2008}] \emph{\bibentry{Fauci2008}}
% \end{labeling}
% }
% \end{multicols}
% \end{WidePar}

\ChapterEnd